\documentclass[a4paper]{article}

\usepackage{amsfonts, amsmath, amssymb, amsthm}
\usepackage{booktabs}
\usepackage{fullpage}
\usepackage{hyperref}

\newtheorem{theorem}{Theorem}[section]
\newtheorem{proposition}[theorem]{Proposition}
\newtheorem{corollary}[theorem]{Corollary}
\newtheorem{conjecture}[theorem]{Conjecture}
\newtheorem{problem}[theorem]{Problem}
\newtheorem{remark}[theorem]{Remark}
\newtheorem{definition}[theorem]{Definition}

\newcommand{\F}{\mathbb{F}}
\newcommand{\N}{\mathbb{N}}
\newcommand{\Z}{\mathbb{Z}}
\newcommand{\abs}[1]{\left|#1\right|}
\DeclareMathOperator{\MCP}{MCP}
\renewcommand{\qedsymbol}{$\blacksquare$}

\setlength{\parindent}{0em}
\setlength{\parskip}{1em}

\begin{document}

\title{Open Problems and Resolved Conjectures for \textit{Lights Out} Nullity}
\author{William Boyles}
\date{\today}
\maketitle

\begin{abstract}
	Let $d(n)$ be the nullity over $\F_2$ of the bounded $n\times n$
	\textit{Lights Out} grid.  Earlier notes in this directory mixed genuine
	open problems with statements subsequently proved in other manuscripts,
	standard Fibonacci-polynomial identities, computational observations,
	and several false parameterizations.  This document is the canonical
	status list as of September 2026.

	The surviving problems concern prime-power stabilization at unknown
	base-$2$ Wieferich primes, prime ``mad'' numbers, the global
	$4/5$ maximum-nullity bound and its record holders, independent
	impossible-nullity seeds, and the Most Clicks Problem beyond nullity two.
	We also record the exact fixed-odd-part theorem that supersedes the old
	$b$-factoring notebook, the computer-assisted proof that $34$ is
	impossible, and explicit counterexamples to the retired formulations.
\end{abstract}

\tableofcontents

\section{Definitions and status conventions}

Let
\[
	F_0(x)=0,\qquad F_1(x)=1,\qquad
	F_n(x)=xF_{n-1}(x)+F_{n-2}(x)
\]
in $\F_2[x]$.  Sutner's formula is
\begin{equation}\label{eq:sutner}
	d(n)=\deg\gcd\!\left(F_{n+1}(x),F_{n+1}(x+1)\right)
\end{equation}
\cite{Sutner2000}.

For a graph $G$, the Most Clicks Problem asks for the largest, over solvable
configurations, of the minimum number of clicks needed to solve the
configuration.  We write $\MCP(n)$ for the answer on the $n\times n$ grid.

The status descriptions below distinguish:
\begin{itemize}
	\item \emph{published theorem}: available in the cited literature;
	\item \emph{repository theorem}: proved in another manuscript in this
	      repository, but not necessarily published;
	\item \emph{computer-assisted repository theorem}: reduced to a finite
	      exact computation whose code and certificate are in the repository;
	\item \emph{open}: no proof or counterexample is known;
	\item \emph{retired}: proved, superseded, malformed, or false, and
	      therefore no longer stated as a conjecture.
\end{itemize}

\section{The fixed-odd-part theorem}
\label{sec:fixed-odd}

This theorem disposes of the old $b$-specific polynomial conjectures and
the meta-conjecture about the recurrence in $k$.

\begin{theorem}[Fixed odd part]\label{thm:fixed-odd}
	Let $b$ be odd, $k\geq1$, $s=k-1$, and
	$\epsilon_b=[\,3\mid b\,]$.  Define
	\[
		C_{b,k}(x)
		=\gcd\!\left(F_{b2^s}(x),F_{b2^s}(x+1)\right).
	\]
	Then
	\begin{equation}\label{eq:gcd-fixed-odd}
		C_{b,k}(x)
		=[x(x+1)]^{\epsilon_b(2^s-1)}
		 C_{b,1}(x)^{2^s}.
	\end{equation}
	Consequently
	\begin{equation}\label{eq:d-fixed-odd}
		d\!\left(b2^{k-1}-1\right)
		=2^s d(b-1)+2\epsilon_b(2^s-1),
	\end{equation}
	and
	\begin{equation}\label{eq:d-fixed-recurrence}
		d\!\left(b2^k-1\right)
		=2d\!\left(b2^{k-1}-1\right)+2[\,3\mid b\,].
	\end{equation}
\end{theorem}

\begin{proof}
	Hunziker, Machiavelo, and Park prove
	\[
		F_{b2^s}(x)=x^{2^s-1}F_b(x)^{2^s}
	\]
	\cite{HunzikerMachiaveloPark2004}; the same identity has recently been
	revisited in the classification of powerful Fibonacci polynomials
	\cite{BatesEtAl2027}.  Substitute $x+1$ for $x$.

	For odd $b$, $F_b$ is the square of a square-free polynomial.
	The factor $x+1$ divides $F_b(x)$ exactly when $3\mid b$, and then its
	multiplicity is exactly two; symmetrically, $x$ divides $F_b(x+1)$
	exactly in the same case.  Neither $x$ nor $x+1$ divides
	$C_{b,1}$.  Taking the gcd of the two displayed factorizations gives
	\eqref{eq:gcd-fixed-odd}.  Degrees give \eqref{eq:d-fixed-odd}, and
	comparing $s$ with $s+1$ gives \eqref{eq:d-fixed-recurrence}.
	The corresponding nullity recurrence is also proved in
	\cite{BoylesSutner}.
\end{proof}

\begin{corollary}\label{cor:31-33}
	For every $k\geq1$,
	\[
		d\!\left(31\cdot2^{k-1}-1\right)=5\cdot2^{k+1},
		\qquad
		d\!\left(33\cdot2^{k-1}-1\right)=11\cdot2^k-2.
	\]
\end{corollary}

\begin{proof}
	The exact evaluations $d(30)=20$ and $d(32)=20$ give the result from
	\eqref{eq:d-fixed-odd}.  Equivalently, the Kloosterman calculation in
	\path{tex/upper_bound/upper_bound.tex} gives $G(31)=20$ and
	$G(33)=22$.
\end{proof}

\begin{remark}
	The explicit formulas formerly conjectured for $b=1,3,5,7,9$ are all
	instances of the published Fibonacci factorization used above.
	The ``zero once, zero always'' conjecture and the claimed recurrence
	$d(b,k+1)\in\{2d(b,k),2d(b,k)+2\}$ are both immediate consequences of
	\eqref{eq:d-fixed-recurrence}.
\end{remark}

\section{Resolved and retired claims}

\begin{center}
\begin{tabular}{p{0.31\textwidth}p{0.18\textwidth}p{0.43\textwidth}}
\toprule
Former claim & Status & Current source or correction\\
\midrule
Rule-90 and rule-150 coefficient recurrences &
Published theorem / direct substitution &
Sutner's coefficient formula and the substitution $x\mapsto x+1$ give the
two cellular-automaton recurrences \cite{Sutner2000}.  They are expository
identities, not conjectures.\\
\addlinespace
$F_{b2^s}=x^{2^s-1}F_b^{2^s}$ and the $b$-factoring iteration &
Published theorem &
Theorem~\ref{thm:fixed-odd}; see
\cite{HunzikerMachiaveloPark2004,BatesEtAl2027}.\\
\addlinespace
The $b=31$ and $b=33$ nullity formulas &
Repository theorem &
Corollary~\ref{cor:31-33} and
\path{tex/upper_bound/upper_bound.tex}.\\
\addlinespace
Eventual constancy of $d(a^k-1)$ for fixed odd $a$ &
Repository theorem &
The finite-support truncation theorem in
\path{tex/multiprime_stabilization/multiprime_stabilization.tex}.\\
\addlinespace
$26$ is impossible, and impossibility propagates by $A\mapsto2A+2$ &
Repository theorem &
\path{tex/finite_fields/finite_fields.tex}.\\
\addlinespace
$34$ is impossible &
Computer-assisted repository theorem &
\path{tex/finite_witness_bounds/finite_witness_bounds.tex}; the independent
exact rerun checks all $32\,767$ candidates and reproduces, for the
$16\,100$ newly computed pairs, SHA-256
\texttt{64cc5c41bf1bf197c833f695e3aaf3308b4154b358667f1a4f79b40a0a2458ab}.\\
\addlinespace
All nullity-$12$ grids have size $340k-256$ &
False &
$d(594)=12$, and $594=340\cdot2-86$.  The corrected necessary classes are
$n\equiv84\pmod{170}$.\\
\addlinespace
The displayed higher-nullity MCP quadratics in the old summary &
False as written &
They used the subsequence index in place of the multiplier
$q=(n+1)/a$; for example the old formula predicts $\MCP(14)=48$, while
the exact value is $123$.\\
\addlinespace
Every varying polynomial gcd is a geometric series with a power-of-two
number of terms &
False &
For $b=15$, $C_{15,1}=1+x^2+x^4$ has three terms, while
$C_{15,2}=x+x^2+x^5+x^6+x^9+x^{10}$ has six and is not a geometric
series.\\
\bottomrule
\end{tabular}
\end{center}

\section{Prime-power stabilization}

For a prime $p>5$, put
\[
	c_p=v_p\!\left(2^{\operatorname{ord}_p(2)}-1\right).
\]
Every non-Wieferich prime has $c_p=1$.

\begin{conjecture}[Wieferich endpoint]\label{conj:prime-power}
	For every prime $p>5$ and every $k\geq1$,
	\[
		d(p^k-1)=d(p-1).
	\]
\end{conjecture}

The statement is a theorem for every non-Wieferich prime.  More generally,
\[
	d(p^k-1)=d(p^{c_p}-1)\qquad(k\geq c_p).
\]
The two known base-$2$ Wieferich primes, $1093$ and $3511$, both have
$c_p=2$, and exact endpoint computations give
\[
	d(p^2-1)=d(p-1)=0.
\]
Thus Conjecture~\ref{conj:prime-power} holds for every currently known
prime and, in particular, for every $p<2^{64}$.

The minimal unresolved assertion is the endpoint lemma: if $\theta$ has
order $p^{c_p}$ and
\[
	1+\theta^a+\theta^{-a}+\theta^b+\theta^{-b}=0,
\]
then $p^{c_p-1}\mid a,b$.  See
\path{tex/five_term_plateau/five_term_plateau.tex} for the equivalent
polynomial, cyclotomic, coding, dynamical, elliptic, and group-theoretic
forms.

\section{Prime mad numbers}

For a prime $p>3$, let
\[
	R_p=
	\left\{\zeta+\zeta^{-1}:
	\operatorname{ord}(\zeta)=p\right\}
	\subseteq\overline{\F}_2.
\]

\begin{proposition}\label{prop:prime-mad}
	For every prime $p>3$,
	\[
		d(p-1)=2\abs{R_p\cap(R_p+1)}.
	\]
	Thus $d(p-1)>0$ exactly when there are nontrivial $p$-th roots
	$t,s$ satisfying
	\[
		1+t+t^{-1}+s+s^{-1}=0.
	\]
\end{proposition}

This is the bounded-grid form of the prime Button Madness problem studied
through Dickson polynomials in \cite{BlokhuisEtAl2018}.

\begin{problem}[Arithmetic characterization]\label{prob:prime-mad}
	Characterize the primes $p$ for which
	$R_p\cap(R_p+1)\neq\varnothing$.
\end{problem}

\begin{conjecture}[Prime infinitude]\label{conj:prime-mad}
	There are infinitely many primes $p$ for which $d(p-1)>0$.
\end{conjecture}

Every Mersenne prime greater than $7$ and every Fermat prime greater than
$3$ has this property; exact Kloosterman formulas appear in
\path{tex/five_term_plateau/five_term_plateau.tex}.  Brouwer and Blokhuis
prove infinitely many mad \emph{integers}, but that does not prove
Conjecture~\ref{conj:prime-mad}, since infinitude of the relevant prime
subfamily remains open \cite{BrouwerMadness}.  The necessary condition
\[
	\rho_2(p)=\min\{r>0:2^r\equiv\pm1\pmod p\}\leq\sqrt p
\]
is not sufficient: $\rho_2(241)=12<\sqrt{241}$ but $d(240)=0$.

\section{Maximum nullity}

\begin{conjecture}[The $4/5$ bound]\label{conj:four-fifths}
	For every $n\geq1$,
	\[
		5d(n)\leq4(n+1),
	\]
	with equality if and only if
	$n=5\cdot2^k-1$ for some $k\geq0$.
\end{conjecture}

The equality family itself is proved:
\[
	d(5\cdot2^k-1)=4\cdot2^k.
\]
The manuscript \path{tex/upper_bound/upper_bound.tex} reduces the global
bound to a multiplicative-subgroup point count on an elliptic curve and
proves it:
\begin{itemize}
	\item when the odd part of $n+1$ is $2^r\pm1$;
	\item under an explicit Weil-type criterion;
	\item for every odd base $b\leq10^5$ by exact computation.
\end{itemize}
In all these regimes $b=5$ is the unique equality case.  No published or
repository proof of the uniform bound for every $b$ is known.
The strongest general published comparison remains Sarkar's theorem
$d(n)=n$ if and only if $n=4$ \cite{Sarkar1996}.

Define the positive record-holder set
\[
	\mathcal M=
	\{n\geq1:d(n)>0\text{ and }d(n)>d(m)\text{ for every }m<n\}.
\]

\begin{conjecture}[Record holders]\label{conj:records}
	\[
		\mathcal M=
		\{5\cdot2^{k-1}-1:k\geq1\}
		\cup\{31\cdot2^{k-1}-1:k\geq1\}
		\cup\{33\cdot2^{k-1}-1:k\geq2\}.
	\]
\end{conjecture}

This is verified for all $n\leq100000$.  The exclusion of
$33\cdot2^0-1=32$ is necessary because $d(32)=20$ only ties the earlier
record at $n=30$.

\section{Impossible nullities}

Let
\[
	I=\{A\in2\N:A\neq d(n)\text{ for every }n\}.
\]
If $A\in I$, call it a \emph{seed} when it is not obtained from a smaller
member of $I$ by iterating $A\mapsto2A+2$.

The repository proves that $26$ and $34$ are seeds.  The doubling
recurrence then gives the two infinite impossible families
\[
	2^t(26+2)-2
	\quad\text{and}\quad
	2^t(34+2)-2
	\qquad(t\geq0).
\]
The finite-witness theorem in
\path{tex/finite_witness_bounds/finite_witness_bounds.tex} further proves
that membership in the image of $d$ is decidable by a finite computation
for each prescribed value.

\begin{conjecture}[Infinitely many seeds]\label{conj:seeds}
	The set $I$ contains infinitely many seed values.
\end{conjecture}

\section{Small nullities and the Most Clicks Problem}

The two-adic formula \eqref{eq:d-fixed-odd}, together with the ranks of the
small translation-invariant Fibonacci factors, gives the following global
necessary conditions.

\begin{theorem}[Small-nullity congruences]\label{thm:small-nullities}
	\[
	\begin{array}{rcl}
	d(n)=4 &\Longrightarrow& n\equiv4\pmod{10},\\
	d(n)=6 &\Longrightarrow& n\equiv11\pmod{24},\\
	d(n)=8 &\Longrightarrow&
		n\equiv9\pmod{20}\ \text{or}\ n\equiv16\pmod{34},\\
	d(n)=10&\Longrightarrow& n\equiv29\pmod{60},\\
	d(n)=12&\Longrightarrow& n\equiv84\pmod{170}.
	\end{array}
	\]
\end{theorem}

\begin{proof}
	Write $n+1=2^sb$, with $b$ odd, and put $D=d(b-1)$.
	For even boards $b-1$, one has $4\mid D$.  Solving
	\[
		d(n)=2^sD+2[\,3\mid b\,](2^s-1)
	\]
	for the five target values gives the possible pairs $(s,D)$.

	If $D=4$, the square-free half of the common Fibonacci factor is
	$x^2+x+1$, of Fibonacci rank $5$, so $5\mid b$.
	If $D=8$, the corresponding translation-invariant factor is
	$x^4+x+1$, of rank $17$, so $17\mid b$.
	If $D=12$, the only viable degree-six factor has rank $85$.
	The two alternative rank-$63$ factors occur together and would force
	nullity at least $24$, so $85\mid b$.  Substitution gives the displayed
	congruences.  The exact rank calculation is recorded in
	\path{tex/finite_witness_bounds/finite_witness_bounds.tex} and the
	supporting nullity code.
\end{proof}

For nullity two, the repository proves that every grid has size $6k-1$ and
\[
	\MCP(6k-1)=26k^2-12k+1;
\]
see \path{tex/nullity2/nullity2.tex} and \cite{BoylesMostClicks}.
No higher-nullity family has been proved.

The following corrected fits use the natural multiplier
$q=(n+1)/a$, not the subsequence index that appeared in the old summary.

\begin{conjecture}[Empirical MCP fits]\label{conj:mcp-fits}
	For each row below, if $n=aq-1$ with $q$ odd and $d(n)$ has the stated
	value, then $\MCP(n)=P(q)$.
	\[
	\begin{array}{ccl}
	d(n)&a&P(q)\\
	\hline
	4&5&17q^2-10q\\
	6&12&88q^2-24q+1\\
	8&10&60q^2-20q-3\\
	8&17&161q^2-34q-3\\
	10&30&506q^2-60q-3
	\end{array}
	\]
\end{conjecture}

The rows agree with respectively $19,6,7,6,$ and $4$ exact values in
\path{code/serialization/min_clicks_answers.txt}, but no proof is known.
The former nullity-$12$ fit is not retained: even $\MCP(84)$ is known only
within the interval $3579\leq\MCP(84)\leq3591$, and the old source used
inconsistent coefficients.

\begin{problem}[Higher-nullity MCP]\label{prob:higher-mcp}
	Determine $\MCP(n)$ on any infinite family of grids with nullity at
	least four.  More generally, determine whether the MCP is polynomial or
	quasi-polynomial in the multiplier $q$ on a fixed nullity family.
\end{problem}

\section{Computational evidence and maintenance}

The nullity data in \path{code/serialization/b159257.txt} contains every
$d(n)$ for $1\leq n\leq100000$.  It verifies:
\begin{itemize}
	\item Conjecture~\ref{conj:four-fifths}, with the stated equality cases,
	      throughout that range;
	\item Conjecture~\ref{conj:records} throughout that range;
	\item all congruences in Theorem~\ref{thm:small-nullities};
	\item the necessity of both nullity-$8$ classes (the first omitted old
	      case is $d(69)=8$);
	\item both mod-$340$ subprogressions of the single nullity-$12$ class
	      $n\equiv84\pmod{170}$ (the first omitted old case is
	      $d(594)=12$, with $594\equiv254\pmod{340}$).
\end{itemize}

The old standalone $b$-factoring, maximum-nullity, and rule-90/rule-150
notes have been removed from the conjectures directory.  Their valid
content is either Theorem~\ref{thm:fixed-odd}, a cited published identity,
or material in the dedicated manuscripts named above.  Generated PDFs and
unused figures were removed with them so that this source remains the only
status document in the directory.

\bibliographystyle{amsalpha}
\bibliography{refs.bib}

\end{document}
